\chapter{Related Work}
\label{chap:related-work}

This chapter positions the project with respect to the main scientific and methodological families used in the report: high-fidelity numerical simulation of Rayleigh--Taylor instability, spectral or multi-resolution representations of physical fields, and generative modeling for surrogate simulation. It should be completed with precise bibliographic references before final submission.

\section{High-Fidelity Simulation of Rayleigh--Taylor Instability}

Direct Numerical Simulations are used as reference data because they resolve the relevant spatial and temporal scales of the flow without relying on turbulence closure models. In the context of RTI, they make it possible to observe the transition from early interface perturbations to non-linear bubble, spike, and mixing structures. Their main limitation is computational cost, which restricts the number of available realizations and motivates data-driven surrogate modeling.

DNS has become a central tool for turbulence research because it gives direct
access to the full resolved flow field and makes it possible to analyze
physical mechanisms without introducing an explicit turbulence model. Moin and
Mahesh describe DNS as a reference approach for studying turbulent flows, while
also emphasizing its strong dependence on computational resources and spatial
resolution \citep{MoinMahesh1998}. This limitation is particularly important in
the present work: RTI simulations are rich enough to describe multi-scale
mixing, but expensive enough that only a limited dataset can be generated.
Pope's reference text on turbulent flows also provides the physical background
for interpreting coherent structures, mixing, spectra, and statistical
quantities in turbulent regimes \citep{pope2000turbulent}.

\section{Spectral and Multi-resolution Representations}

Fourier methods are natural for periodic flows and turbulence analysis because they describe a field in terms of global frequency modes. They are particularly useful for energy-spectrum analysis and pseudo-spectral solvers. However, global basis functions are less adapted to localized structures and sharp gradients.

Wavelet methods provide a complementary representation localized in both space and scale. This property is useful for RTI fields, where large coherent structures coexist with localized small-scale details. The multi-resolution structure also makes wavelets attractive for compression, denoising, and scale-aware learning.

Wavelet representations are especially relevant for this project because they
combine spatial localization with scale localization. Mallat's work presents
wavelet decompositions as a multi-resolution signal representation able to
separate coarse structures from localized details \citep{Mallat2009}. This
property is directly connected to RTI fields, where large interface
deformations coexist with sharper gradients and smaller mixing structures. In
contrast with purely Fourier representations, wavelets therefore offer a
natural way to introduce scale-aware information into the generative model.

\section{Generative Surrogate Models for Physical Fields}

Generative models aim to learn a probability distribution from reference samples and then generate new realizations from this distribution. Diffusion models are particularly relevant because they can represent complex high-dimensional distributions and have shown strong empirical performance in image generation. For physical surrogate modeling, the challenge is not only visual realism but also the preservation of statistical and physical properties.

Denoising Diffusion Probabilistic Models provide the main generative baseline
used in this report. Ho, Jain, and Abbeel introduced a formulation in which a
model learns to reverse a gradual noising process, making it possible to sample
complex high-dimensional distributions from noise \citep{Ho2020DDPM}. Song et
al. later connected diffusion and score-based generative modeling through
stochastic differential equations, giving a broader continuous-time view of the
same family of methods \citep{Song2021SDE}. The present work differs from these
general image-generation settings by applying diffusion models to RTI density
fields generated by DNS. The contribution is therefore not only to generate
visually plausible samples, but also to respect the physical symmetries of the
dataset through adapted augmentations, to evaluate generated fields with
physics-oriented metrics such as PhyFID and profile or spectral statistics, and
to explore wavelet conditioning as a scale-aware representation for turbulent
fields.

\section{Motivation for the Chosen Direction}

The reviewed approaches suggest that a direct diffusion baseline is a necessary reference point, while wavelet representations may introduce a useful inductive bias for limited and multi-scale physical datasets. This motivates the two-pipeline strategy developed in the following chapters.
